\scanpage{13}

Donc on a un foncteur

\diagram{\begin{tikzcd}[ampersand replacement=\&,column sep=large]
{\begin{array}{c}\text{Géom. d'incidences}\\\text{(avec iso)}\end{array}}
\arrow[r]\&
{\begin{array}{l}\unsure{G|\mathcal G|\mathfrak G|\mathfrak S}_2\text{-ens. où $\sigma_0,\sigma_1,\sigma_2$}\\
 \text{et $\sigma_0\sigma_2$ opèrent ss pts}\\ \text{fixes (avec iso)}\end{array}}
\end{tikzcd}}

On va voir que \inserted{\uline{Thm}} ce foncteur est \uline{pleinement fidèle} -- donc qu'on peut réconstruire $(S,A,F,R,R')$ à partir de $E=\operatorname{Rep}$ en \deleted{tant} que $G_2$-ens. Notons qu'on a

\[
\begin{aligned}
 E/(\sigma_1,\sigma_2)&\longrightarrow S=\Phi_0\\
 E/(\sigma_0,\sigma_2)&\longrightarrow A=\Phi_1\\
 E/(\sigma_0,\sigma_1)&\longrightarrow F=\Phi_2
\end{aligned}
\]

\deleted{qui sont des applications \uline{surjectives}.}

Il est clair que l'application médiane est un iso [$\sigma_0,\sigma_2$ engendrent un ss-gpe de $G_2$ $\simeq\mathbb Z_2\times\mathbb Z_2$, qui opère librement sur $E$ \deleted{\unsure{par hyp.},} le quotient est $\xrightarrow{\sim} A$]. Or les axiomes 3) et 3$'$) (pas encore utilisés) expriment respectivement \deleted{\unsure{que sont surjectives}} que la dernière et la première flèches sont \deleted{\unsure{aussi}} \uline{bijectives}. Celà permet d'exprimer $S,A,F$ à partir de $R$.

% archive sheet number: 12

